\documentclass[a4paper, 10pt]{article}

\usepackage[a4paper, top=1.2cm, bottom=1.2cm, left=1.4cm, right=1.4cm]{geometry}
\usepackage{fontspec}
\usepackage[ngerman]{babel}
\setmainfont{Libertinus Serif}

\usepackage{xcolor}
\usepackage[most]{tcolorbox}
\usepackage{amsmath, amssymb}
\usepackage{fontawesome5}
\usepackage{enumitem}

% Farben
\definecolor{primaryBlue}{RGB}{20, 50, 100}
\definecolor{lightBlue}{RGB}{240, 245, 253}
\definecolor{lvlGreen}{RGB}{25, 125, 55}
\definecolor{lvlBlue}{RGB}{0, 85, 179}
\definecolor{lvlPurple}{RGB}{115, 30, 135}

\definecolor{step1Bg}{RGB}{255, 253, 240}
\definecolor{step1Border}{RGB}{230, 200, 100}
\definecolor{step1Text}{RGB}{140, 95, 0}

\definecolor{step2Bg}{RGB}{245, 250, 255}
\definecolor{step2Border}{RGB}{150, 200, 245}
\definecolor{step2Text}{RGB}{0, 90, 155}

\definecolor{step3Bg}{RGB}{242, 250, 244}
\definecolor{step3Border}{RGB}{140, 215, 160}
\definecolor{step3Text}{RGB}{20, 115, 50}

\setlength{\parindent}{0pt}
\pagestyle{empty}

\begin{document}

% ==============================================================================
% SEITE 1: TIPP-KARTEN ERARBEITUNG: DIVISION UND BRÜCHE
% ==============================================================================

\noindent
{\color{primaryBlue}\bfseries\large \faLayerGroup\enskip Gestufte Hilfen: Arbeitsblatt „Division und Brüche“}
\hfill
{\footnotesize\color{primaryBlue!80} \textbf{Mathematik 6D} \quad$|$\quad Herr Dayi}
\vspace{1.5mm}
{\color{primaryBlue!40}\hrule height 0.8pt}
\vspace{3mm}

% --- AUFGABE 1 ---
\begin{tcolorbox}[enhanced, colback=lvlBlue!5, colframe=lvlBlue!60, arc=2mm, boxrule=0.8pt, left=3mm, right=3mm, top=2mm, bottom=2mm]
    {\bfseries\color{lvlBlue} \faThLarge\enskip Aufgabe 1 (Kernaufgabe): Kärtchen-Quartette färben}\\[1mm]
    {\footnotesize\textit{Färbe zusammengehörige Kärtchen in der gleichen Farbe ein (Bilder, Divisionen, Brüche, Werte).}}
\end{tcolorbox}

\vspace{1.2mm}

\begin{tcolorbox}[colback=step1Bg, colframe=step1Border, arc=1.8mm, boxrule=0.7pt, left=3mm, right=3mm, top=1.8mm, bottom=1.8mm]
    {\color{step1Text}\small\bfseries\faEye\enskip Stufe 1: Anschauung}\\[1mm]
    {\footnotesize Schau dir die vier Bilder genau an: \textbf{A} (Streifen 300 in 4 Teile), \textbf{C} (20 Murmeln in 5 Röhren), \textbf{F} (Waage: 2 kg auf 4 Schalen) und \textbf{I} (Streifen 2 in 3 Teile).}
\end{tcolorbox}

\begin{tcolorbox}[colback=step2Bg, colframe=step2Border, arc=1.8mm, boxrule=0.7pt, left=3mm, right=3mm, top=1.8mm, bottom=1.8mm]
    {\color{step2Text}\small\bfseries\faLightbulb\enskip Stufe 2: Denkanstoß}\\[1mm]
    {\footnotesize Finde zu jedem Bild die Rechnung ($:$), den Bruch ($\frac{\text{Zähler}}{\text{Nenner}}$) und das Ergebnis:
    z.\,B. Streifen 300 auf 4 Teile $\to 300:4 \to \frac{300}{4} \to 75$.}
\end{tcolorbox}

\begin{tcolorbox}[colback=step3Bg, colframe=step3Border, arc=1.8mm, boxrule=0.7pt, left=3mm, right=3mm, top=1.8mm, bottom=1.8mm]
    {\color{step3Text}\small\bfseries\faTasks\enskip Stufe 3: Vorgehens-Leitfaden \& Kontrolltipp}\\[1mm]
    {\footnotesize
    $\bullet$ Starte mit den ganzen Zahlen: Bild C ($20$ Murmeln in $5$ Röhren $\to$ welche Zahl ist $20 : 5$?).\\[0.5mm]
    $\bullet$ Nimm Bild A ($300$ aufgeteilt in $4$ Abschnitte $\to$ rechne $300 : 4$).\\[0.5mm]
    $\bullet$ Bild I ($2$ Ganze in $3$ Teile $\to 2 : 3 = \frac{2}{3}$) und Waage F ($2$ kg auf $4$ Schalen $\to \frac{2}{4}$).\\[0.5mm]
    \textit{Kontrolltipp:} Genau 3 Kärtchen (Buchstaben M, N, P) bleiben am Ende übrig und gehören in keine Gruppe!
    }
\end{tcolorbox}

\vspace{3.5mm}

% --- AUFGABE 2 ---
\begin{tcolorbox}[enhanced, colback=lvlBlue!5, colframe=lvlBlue!60, arc=2mm, boxrule=0.8pt, left=3mm, right=3mm, top=2mm, bottom=2mm]
    {\bfseries\color{lvlBlue} \faPuzzlePiece\enskip Aufgabe 2 (Kernaufgabe): Merksatz-Puzzle}\\[1mm]
    {\footnotesize\textit{Bringe die Schnipsel in die richtige Reihenfolge und schreibe die beiden Sätze auf.}}
\end{tcolorbox}

\vspace{1.2mm}

\begin{tcolorbox}[colback=step1Bg, colframe=step1Border, arc=1.8mm, boxrule=0.7pt, left=3mm, right=3mm, top=1.8mm, bottom=1.8mm]
    {\color{step1Text}\small\bfseries\faEye\enskip Stufe 1: Anschauung}\\[1mm]
    {\footnotesize Suche nach Satzanfängen (Großbuchstaben: \textit{„Wenn man...“} und \textit{„Auf alle Fälle...“}) sowie nach Satzenden mit Punkt!}
\end{tcolorbox}

\begin{tcolorbox}[colback=step2Bg, colframe=step2Border, arc=1.8mm, boxrule=0.7pt, left=3mm, right=3mm, top=1.8mm, bottom=1.8mm]
    {\color{step2Text}\small\bfseries\faLightbulb\enskip Stufe 2: Denkanstoß}\\[1mm]
    {\footnotesize Satz 1 beschreibt, ob das Ergebnis eine ganze Zahl ist oder nicht ($16:4$ vs. $9:4$). Satz 2 stellt die Brücke zum Bruchstrich her!}
\end{tcolorbox}

\begin{tcolorbox}[colback=step3Bg, colframe=step3Border, arc=1.8mm, boxrule=0.7pt, left=3mm, right=3mm, top=1.8mm, bottom=1.8mm]
    {\color{step3Text}\small\bfseries\faPenNib\enskip Stufe 3: Satzgerüst zur Hilfe}\\[1mm]
    {\footnotesize
    Ordne die Schnipsel so, dass sie diese logische Aussage ergeben:\\[0.5mm]
    \textbf{Satz 1:} „Wenn man zwei natürliche Zahlen dividiert, \dots\dots\dots\dots\dots\dots\dots (z.\,B. bei $16:4$) oder nicht (z.\,B. bei $9:4$).“\\[0.8mm]
    \textbf{Satz 2:} „Auf alle Fälle lässt sich jede Division \dots\dots\dots\dots\dots\dots darstellen. Hier also $\frac{16}{4}$ bzw. $\frac{9}{4}$.“\\[0.5mm]
    \textit{Kontrolltipp:} Achte auf Satzanfänge (Großschreibung) und Punkte am Satzende!
    }
\end{tcolorbox}

\newpage

% ==============================================================================
% SEITE 2: DIFFERENZIERUNG & SPRINTER (AUFGABE 3 & 4)
% ==============================================================================

\noindent
{\color{primaryBlue}\bfseries\large \faRunning\enskip Gestufte Hilfen: Sprinteraufgaben (Aufgabe 3 \& 4)}
\hfill
{\footnotesize\color{primaryBlue!80} \textbf{Mathematik 6D} \quad$|$\quad Herr Dayi}
\vspace{1.5mm}
{\color{primaryBlue!40}\hrule height 0.8pt}
\vspace{3mm}

% --- AUFGABE 3 ---
\begin{tcolorbox}[enhanced, colback=lvlGreen!5, colframe=lvlGreen!60, arc=2mm, boxrule=0.8pt, left=3mm, right=3mm, top=2mm, bottom=2mm]
    {\bfseries\color{lvlGreen} \faStar\enskip Aufgabe 3 (Sprinter): Ist das Ergebnis eine natürliche Zahl?}\\[1mm]
    {\footnotesize\textit{Kreuze an, in welchen Fällen die Division ohne Rest aufgeht.}}
\end{tcolorbox}

\vspace{1.2mm}

\begin{tcolorbox}[colback=step1Bg, colframe=step1Border, arc=1.8mm, boxrule=0.7pt, left=3mm, right=3mm, top=1.8mm, bottom=1.8mm]
    {\color{step1Text}\small\bfseries\faEye\enskip Stufe 1: Anschauung}\\[1mm]
    {\footnotesize Bruchstrich = Geteilt durch ($:$). Frage dich: Lässt sich der Zähler (oben) ohne Rest durch den Nenner (unten) dividieren?}
\end{tcolorbox}

\begin{tcolorbox}[colback=step2Bg, colframe=step2Border, arc=1.8mm, boxrule=0.7pt, left=3mm, right=3mm, top=1.8mm, bottom=1.8mm]
    {\color{step2Text}\small\bfseries\faLightbulb\enskip Stufe 2: Denkanstoß}\\[1mm]
    {\footnotesize Ist der Zähler kleiner als der Nenner (z.\,B. $\frac{11}{22} = \frac{1}{2}$), kann das Ergebnis niemals eine natürliche Zahl ($\ge 1$) sein!}
\end{tcolorbox}

\begin{tcolorbox}[colback=step3Bg, colframe=step3Border, arc=1.8mm, boxrule=0.7pt, left=3mm, right=3mm, top=1.8mm, bottom=1.8mm]
    {\color{step3Text}\small\bfseries\faCheckSquare\enskip Stufe 3: Rechengerüst \& Kontrollhilfe}\\[1mm]
    {\footnotesize
    $\bullet$ Prüfe im Kopf mit der Umkehraufgabe: Ist der Zähler ein Vielfaches des Nenners?\\
    $\to$ Beispiel a) $\frac{14}{2}$: Passt die $2$ ohne Rest in die $14$? ($2 \cdot 7 = 14 \implies$ Ja!)\\[0.5mm]
    $\to$ Beispiel c) $\frac{15}{7}$: Passt die $7$ ohne Rest in die $15$? ($7 \cdot 2 = 14$, Rest $1 \implies$ Nein!)\\[0.5mm]
    \textit{Kontrollhilfe:} Bei genau \textbf{7 der 12 Aufgaben} gehört ein Kreuz!
    }
\end{tcolorbox}

\vspace{3.5mm}

% --- AUFGABE 4 ---
\begin{tcolorbox}[enhanced, colback=lvlPurple!5, colframe=lvlPurple!60, arc=2mm, boxrule=0.8pt, left=3mm, right=3mm, top=2mm, bottom=2mm]
    {\bfseries\color{lvlPurple} \faExclamationTriangle\enskip Aufgabe 4 (Experten): Teilen durch Null?}\\[1mm]
    {\footnotesize\textit{In welchen Fällen kannst du die Division ausführen? Gib das Ergebnis an.}}
\end{tcolorbox}

\vspace{1.2mm}

\begin{tcolorbox}[colback=step1Bg, colframe=step1Border, arc=1.8mm, boxrule=0.7pt, left=3mm, right=3mm, top=1.8mm, bottom=1.8mm]
    {\color{step1Text}\small\bfseries\faEye\enskip Stufe 1: Anschauung}\\[1mm]
    {\footnotesize $\frac{0}{3}$ bedeutet: $0$ Euro werden gerecht an $3$ Kinder verteilt (jeder erhält $0$ Euro). Aber $\frac{3}{0}$ bedeutet: $3$ Euro an $0$ Kinder verteilen – das geht nicht!}
\end{tcolorbox}

\begin{tcolorbox}[colback=step2Bg, colframe=step2Border, arc=1.8mm, boxrule=0.7pt, left=3mm, right=3mm, top=1.8mm, bottom=1.8mm]
    {\color{step2Text}\small\bfseries\faLightbulb\enskip Stufe 2: Denkanstoß}\\[1mm]
    {\footnotesize \textbf{Goldene Mathe-Regel:} Durch Null darf man niemals dividieren! Im Nenner (unten) darf keine $0$ stehen.}
\end{tcolorbox}

\begin{tcolorbox}[colback=step3Bg, colframe=step3Border, arc=1.8mm, boxrule=0.7pt, left=3mm, right=3mm, top=1.8mm, bottom=1.8mm]
    {\color{step3Text}\small\bfseries\faCheckSquare\enskip Stufe 3: Strukturhilfe zur Begründung}\\[1mm]
    {\footnotesize
    $\bullet$ Steht die $0$ \textbf{unten im Nenner}, teilst du durch $0$ $\longrightarrow$ \textbf{Nicht ausführbar} (mathematisch nicht definiert!).\\[0.5mm]
    $\bullet$ Steht die $0$ \textbf{oben im Zähler} und unten eine Zahl $\neq 0$, ist das Teilen erlaubt $\longrightarrow$ \textbf{Ausführbar}! Überlege, welche feste Zahl immer herauskommt, wenn man $0$ Dinge verteilt.\\[0.5mm]
    \textit{Kontrolltipp:} Genau 2 Aufgaben sind ausführbar (Ergebnis ist identisch); 2 Aufgaben sind nicht ausführbar.
    }
\end{tcolorbox}

\end{document}
